
Messenger RNA is an tractable drug target due to its accessibility with nucleic acids. Nucleic acid drugs target messenger RNA upstream before it is translated into a functional protein and act by blocking expression through transcriptional interference or gene editing~\cite{Crooke2021, Rinaldi2018}. Antisense oligonucleotides (ASOs) are particularly flexible in terms of the mechanism of action based upon chemical composition and design. They can trigger cleavage of the target RNA by RNase H, prevent translation of the target RNA into protein, or correct splicing errors~\cite{Bennett2019}. As there are now many clinically-approved ASO therapeutics and numerous additional candidates currently being evaluated in clinical trials, RNA-targeted therapy has become a significant tool in therapeutic development~\cite{Crooke2021}.

However, translating an effective ASO requires a specific RNA target sequence that continues to prove challenging. The major challenge in developing an effective ASO is structural. Messenger RNA does not exist as a linear strand of unstructured bases; rather, it folds into secondary structures (such as hairpin loops), forms tertiary contact interactions, and dynamically samples multiple conformations throughout time. Therefore, simply analyzing the potential target sequence will not give you information about which parts of the RNA target will be available for interaction with an antisense oligonucleotide~\cite{Watts2009, Liang2011}; a site predicted to be single-stranded may still be physically inaccessible due to its location within the overall tertiary structure or be transiently blocked by changing stems during sampling of different conformations. Because ASO design depends on both sequence-specific complementary base pairing and physical accessibility to the target site, designing an effective ASO requires a combination of structural and sequence-based knowledge.

Computational methods offer a natural approach to addressing the challenges of developing effective ASOs, although current methods provide limited solutions. Tools for predicting the secondary structure of RNA (e.g., ViennaRNA) predict minimum-energy RNA structures and have been widely used to select ASO target sites~\cite{Lorenz2011, Mathews2004}. Advances in predicting tertiary structure have also emerged (e.g., NuFold) providing new insights into the overall three-dimensional structure of RNA~\cite{Kagaya2025}. However, despite these advances, current methods are fundamentally static -- they can describe the expected structures formed by the RNA under conditions of equilibrium, but they fail to describe how accessability varies under nonequilibrium conditions (i.e., thermal fluctuations), how long a nucleotide remains exposed to solvents, and/or how tightly bound an ASO is to a particular site once it arrives at that site. Molecular dynamics simulation represents one approach toward describing these processes, however, performing molecular-dynamics simulations of systems large enough to model entire proteins is computationally expensive. Coarse-graining the system provides a cost-effective method for simulating larger systems; however, coarse-grained RNA models have not been extensively utilized for identifying candidate ASO target-sites nor for evaluating/ranking target sites.
This thesis develops a computational pipeline capable of identifying candidate sites for ASO binding and evaluating them in terms of both structural accessibility and binding affinity. The starting point for this process is either experimentally-derived PDB coordinates or tertiary structure predictions from NuFold~\cite{Kagaya2025}. Each sequence is subsequently simulated in LAMMPS~\cite{Thompson2022} utilizing a single-interaction-site (SIS) coarse-grained RNA force-field~\cite{Aierken2024}. Target sites identified by secondary-structure analysis via ViennaRNA and/or by per-nucleotide solvent-accessible surface area calculations are subsequently evaluated/ranked using a trajectory-based binding framework employing center-of-mass distance measures similar to those employed in previous protein-RNA binding studies~\cite{Alston2023, Tesei2021}. The final output is a set of quantitative metrics describing the relative binding affinity of each candidate site to an ASO directly from simulation data. The resulting methodology has been tested over a variety of benchmark RNA targets and longer target sequences designed specifically for testing purposes.
